\subsection{代表タスクのケース分析}\label{sec5.cases}
定量評価で見られる失敗とタスクに必要な判断を具体化するため，構成要素比較と同じ状態変化2件・配置2件を対象に，初期状態，ゴール，参照手順，実行過程を分析する．著者による説明的なケース分析である．以下はWindows VHの固定位置条件におけるケース分析用試行のログに基づき，構成要素比較の240件には合算しない．

状態変化\texttt{test\_task1}は，4個のlightswitch（71，173，261，427）が初期状態ですべてONであり，参照列長は0である．必要なのは指示の動詞に応じて操作を始めることではなく，対象全体の初期充足を確認して終了する判断である．

同じ指示の\texttt{test\_task6}では，173と427だけがOFFである．参照手順は173への移動とON操作，427への移動とON操作の4行動であり，ON済みの71と261との区別が必要となる．固定位置試行のGPT-4o mini・記号検証条件では，173への移動とON操作，261への移動が成功した後，427がOFFのまま終了判定がEndとなった．各行動の前提条件を満たすことと，タスク全体を達成することは異なる．この例では検証による再生成は起きていなかった．

配置\texttt{test\_task1}は，cupcake 195と196をkitchencounter 238に置く課題である．参照列は各物体への移動，把持，配置先への移動，配置を繰り返す8行動である．対象を2個とも追跡するとともに，保持・近接と物体の操作適合性を区別する必要がある．参照列の再生では最初の3行動が成功したが，4行動目のPUTが失敗した．保持・近接を確認できても実行可能性の十分条件とはならず，一般的なエラー応答だけから物理的原因を特定することはできなかった．またLLM検証条件では，OPEN deskという候補を棄却した後にOPEN cupcakeへ再生成し，VHが失敗を返す例があった．再生成後に再検証しない実装上の限界も確認できる．

配置\texttt{test\_task65}は，plum 53と54を初期CLOSEDのfridge 103に入れる課題である．参照列は1個目の移動・把持・冷蔵庫への移動・開扉・格納と，2個目の移動・把持・移動・格納の9行動であり，再生に成功した．途中の保持物と残り対象を管理する必要がある．検証なしのGPT-4o miniでは，53を格納し，54を右手に持ったままEndを返す例があった．保持情報は終了判定の入力にも含まれていたため，この失敗を情報抽出漏れだけでは説明できない．なお，両物体のINSIDE関係がゴールであり，最後に冷蔵庫を閉めることを評価時に追加要求してはいない．

これらは失敗過程と必要な判断を示す例である．4ケースの説明内容は2026年9月6日に確認を受けた．確認範囲は，本節の4ケースの初期条件・参照手順・必要な判断・難しい点の説明であり，確認者自身による実行や所要時間の計測は含まない．

行動列の実行可能性とゴール充足の関係を調べるため，基本性能評価の同じ4タスク・2モデルから記録された行動列8本を固定位置で再生した．8本とも記録列全体を再生でき，最終ゴールを満たしたのは5本だった．GPT-4o・Multiの配置\texttt{test\_task65}では10行動をすべて再生し，ゴールを満たした．これは\ref{sec5.replay}の棄却候補8件とは分析対象が異なる．記録されていない終端の失敗行動は含まず，生成時の人物位置との一致も確認していないため，記録接頭列を固定位置で実行した結果として解釈する．各列の長さとゴール充足は表\ref{tab:saved_replay_internal}に示す．
